\section{The table of liturgical days}
\label{sec:table}

\subsection{Its authority}

\begin{normtext}{\loc{nualc} 59}
Præcedentia inter dies liturgicos, quoad eorum celebrationem, unice regitur sequenti tabula.
\end{normtext}

\latin{Unice} --- solely. Precedence among liturgical days is governed by this table and by nothing
else. That is a strong claim and it is meant strictly: there is no residual customary precedence, no
appeal to the dignity of a titular, no general principle that a solemnity always wins. Where two
celebrations fall on one day, the question is only which of them stands higher in the thirteen
numbered places below.

The table is printed immediately after \loc{nualc}~59 and before the General Roman Calendar. It has
no article number of its own; citations to it are conventionally by its three Roman-numeral parts
and the Arabic number of the place within the whole sequence, which runs 1 to 13 continuously across
the parts. This reference follows that convention: ``table n.~8\,b'' means the second entry under
place~8, in part~II.

\subsection{The table in full}

The Latin is given as printed in the 2002 \latin{editio typica tertia}; the right-hand column is an
editorial gloss identifying what each place contains and, where useful, why it sits where it does.
The gloss is not a translation offered for liturgical use.

\begin{latintable}
\textbf{Part I} & Places 1--4: the Triduum, the greatest days of the Lord and their protected neighbours, the solemnities of the General Roman Calendar, and proper solemnities.\\
\textbf{1.} Triduum paschale Passionis et Resurrectionis Domini. & The Paschal Triduum, alone at the
head. \loc{nualc}~18 calls it the summit of the whole liturgical year.\\
\textbf{2.} Nativitas Domini, Epiphania, Ascensio et Pentecostes. Dominicæ Adventus, Quadragesimæ et
Paschæ. Feria IV Cinerum. Feriæ Hebdomadæ sanctæ, a feria II ad feriam V inclusive. Dies infra
octavam Paschæ. & Four solemnities of the Lord, the Sundays of the three strong seasons, Ash
Wednesday, Monday to Thursday of Holy Week, and the weekdays of the Octave of Easter. Note that
Ash Wednesday and the Holy Week weekdays outrank every solemnity, and that a weekday of the Easter
octave does the same.\\
\textbf{3.} Sollemnitates Domini, beatæ Mariæ Virginis, et Sanctorum in Calendario generali
inscriptæ. Commemoratio omnium fidelium defunctorum. & Every other solemnity of the General Roman
Calendar, together with the Commemoration of All the Faithful Departed on 2 November, which is
placed here rather than among feasts.\\
\textbf{4.} Sollemnitates propriæ, nempe: a) Sollemnitas Patroni principalis loci seu oppidi aut
civitatis. b) Sollemnitas dedicationis et anniversarii dedicationis ecclesiæ propriæ.
c) Sollemnitas Tituli ecclesiæ propriæ. d) Sollemnitas aut Tituli, aut Fundatoris, aut Patroni
principalis Ordinis seu Congregationis. & Proper solemnities: the principal patron of the place,
town or city; the dedication and its anniversary for the church itself; the title of the church
itself; and, for religious, the title, founder or principal patron of the order or congregation.
These are the only celebrations in part~I that a General Roman Calendar cannot supply.\\
\textbf{Part II} & Places 5--9: feasts of the Lord, the Sundays of Christmas Time and Ordinary Time, other feasts of the general calendar, proper feasts, and the privileged weekdays.\\
\textbf{5.} Festa Domini in Calendario generali inscripta. & Feasts of the Lord in the General Roman
Calendar. These, and only these, displace an ordinary Sunday under \loc{nualc}~5.\\
\textbf{6.} Dominicæ temporis Nativitatis et dominicæ ``per annum''. & The Sundays of Christmas Time
and of Ordinary Time --- below feasts of the Lord and above every other feast.\\
\textbf{7.} Festa beatæ Mariæ Virginis et Sanctorum Calendarii generalis. & Feasts of the Blessed
Virgin and of the saints in the General Roman Calendar. Being below place~6, such a feast is simply
omitted in a year in which it falls on a Sunday.\\
\textbf{8.} Festa propria, nempe: a) Festum Patroni principalis diœcesis. b) Festum anniversarii
dedicationis ecclesiæ cathedralis. c) Festum Patroni principalis regionis aut provinciæ, nationis,
amplioris territorii. d) Festum Tituli, Fundatoris, Patroni principalis Ordinis seu Congregationis
et provinciæ religiosæ, salvis præscriptis sub n.~4. e) Alia festa alicui ecclesiæ propria.
f) Alia festa inscripta in calendario cuiusque diœcesis vel Ordinis seu Congregationis. & Proper
feasts, in a fixed internal order: diocesan principal patron, cathedral dedication anniversary,
patron of region or province or nation or wider territory, religious title or founder or patron,
other feasts proper to a particular church, other feasts in a diocesan or religious calendar.\\
\textbf{9.} Feriæ Adventus a die 17 ad 24 decembris inclusive. Dies infra octavam Nativitatis.
Feriæ Quadragesimæ. & The privileged weekdays. They stand above every memorial but below every
feast.\\
\textbf{Part III} & Places 10--13: obligatory memorials of the general calendar, proper obligatory memorials, optional memorials, and the ordinary weekdays.\\
\textbf{10.} Memoriæ obligatoriæ Calendarii generalis. & Obligatory memorials of the General Roman
Calendar.\\
\textbf{11.} Memoriæ obligatoriæ propriæ, nempe: a) Memoriæ Patroni secundarii loci, diœcesis,
regionis aut provinciæ religiosæ. b) Aliæ memoriæ obligatoriæ inscriptæ in calendario cuiusque
diœcesis, vel Ordinis seu Congregationis. & Proper obligatory memorials: secondary patrons, and
other obligatory memorials of a diocesan or religious calendar.\\
\textbf{12.} Memoriæ ad libitum, quæ tamen, modo quidem peculiari in Institutionibus generalibus
Missalis Romani et de Liturgia Horarum descripto, fieri possunt etiam diebus de quibus sub n.~9.
Hac eadem ratione, ut memoriæ ad libitum celebrari possunt memoriæ obligatoriæ, quæ accidentaliter
occurrunt in feriis Quadragesimæ. & Optional memorials --- which, in the particular manner described
in the two General Instructions, may be made even on the privileged weekdays of place~9. By the same
provision, obligatory memorials that happen to fall on weekdays of Lent may be kept as optional
memorials.\\
\textbf{13.} Feriæ Adventus usque ad diem 16 decembris inclusive. Feriæ temporis Nativitatis a die 2
ianuarii ad sabbatum post Epiphaniam. Feriæ temporis paschalis a feria II post octavam Paschæ ad
sabbatum ante Pentecosten inclusive. Feriæ ``per annum''. & The ordinary weekdays, in four named
runs. Every other day of the year has appeared above.\\
\end{latintable}

\subsection{Reading the table correctly}

\subsubsection{The three parts are not three grades}

It is tempting to read part~I as ``solemnities,'' part~II as ``feasts'' and part~III as
``memorials.'' The table does not divide that way. Ash Wednesday and the weekdays of Holy Week sit
in part~I though they are weekdays; the Sundays of Ordinary Time sit in part~II between two ranks of
feast; the privileged weekdays of Advent and Lent sit at the bottom of part~II above every memorial.
The parts are a rough grouping for the eye; the operative unit is the numbered place.

\subsubsection{Place 12 contains a permission, not merely a rank}

Place~12 is the only entry in the table that carries a substantive rule inside it. It says that
optional memorials may be made even on the days of place~9 --- the privileged weekdays --- in the
particular manner described in the General Instructions. That manner is set out at \loc{igmr}~355\,a:
on the weekdays of Advent from 17 December, within the Octave of the Nativity, and on the weekdays
of Lent, the Mass is of the occurring liturgical day, but the collect may be taken from a memorial
inscribed that day in the general calendar, provided the day is not Ash Wednesday or a weekday of
Holy Week. So a memorial that has lost on the table is not simply erased; it survives as an oration.
This is the postconciliar residue of what older systems called a commemoration, and it is the reason
a precedence table alone never settles what is actually said at an altar.

\subsubsection{The table settles occurrence; it does not settle the choice of Mass}

\loc{nualc}~59 governs \latin{præcedentia inter dies liturgicos}. Which Mass a particular priest may
say on a day whose identity is thereby settled is a different question, governed by
\loc{igmr}~353--355. On solemnities he is bound to the calendar of the church in which he celebrates
(353). On Sundays, on the weekdays of Advent, Christmas Time, Lent and Easter Time, and on feasts
and obligatory memorials, a Mass with a congregation follows the calendar of the church of
celebration, while a Mass at which a single minister participates may follow either that calendar or
his own (354). On optional memorials the choice is wide, and on weekdays in Ordinary Time it is at
its widest: the Mass of the weekday, of an occurring optional memorial, of any saint inscribed that
day in the Martyrology, or a Mass for various needs or a votive Mass (355\,c). \loc{igmr}~355 adds
two cautions that belong here rather than in a rubric book: the priest is not to omit the assigned
weekday readings frequently and without sufficient cause, because the Church desires that a richer
table of God's word be spread for the faithful; and where a choice lies between a memorial of the
general calendar and one of a diocesan or religious calendar, other things being equal and following
tradition, the particular memorial is to be preferred.

\subsubsection{Sunday's own protection is stated twice}

\loc{nualc}~5 and table places 2, 5 and 6 say the same thing in two registers. A Sunday of Advent,
Lent or Easter is in place~2 and cannot be displaced at all. A Sunday of Christmas Time or Ordinary
Time is in place~6 and is displaced only by places 1--5: that is, by the Triduum, by the four great
solemnities and their neighbours in place~2, by any solemnity of the General Roman Calendar
(place~3), by a proper solemnity (place~4), or by a feast of the Lord (place~5). Nothing else
reaches it. There is one further permission, and it is a permission rather than a precedence rule:
\loc{nualc}~58 allows, for the pastoral good of the faithful, that on Sundays in Ordinary Time
celebrations occurring during the week which are dear to the piety of the faithful may be observed,
\emph{provided} those celebrations rank above the Sunday in the table of precedence. Since a
celebration above place~6 would displace the Sunday anyway, what \loc{nualc}~58 actually authorises
is the transfer of such a celebration \emph{to} the Sunday, and it authorises it for all Masses
celebrated with a congregation.
